\section{Critical Timeline}

This chronology records documentary, juridical, devotional, and biographical events without assigning them one authority.

\begingroup
\small
\begin{description}[style=nextline,leftmargin=1.5em,labelsep=0.35em,itemsep=0.35em]
  \item[7 November 1831] Françoise Mélanie Mathieu Calvat is born at Corps.
  \item[August 1835] Pierre Maximin Giraud is born at Corps; civil-date references vary between 26 and 27 August.
  \item[1845--1847] Potato disease, grain distress, high prices, and famine conditions affect France and Europe with substantial regional differences.
  \item[19 September 1846] Maximin, eleven, and Mélanie, fourteen, report the encounter with the luminous weeping Beautiful Lady on Mont Planeau. The public discourse, two separate private communications, and final commission belong to the reported event.
  \item[20 September 1846] The Pra relation, the first written account, is composed in Mélanie's presence. The lost original survives through Abbé Lagier's certified February 1847 copy and early derivatives.
  \item[October 1846] Early documents attest the children's claim that each received a separate private communication.
  \item[February 1847] Abbé Lagier copies the Pra relation and records early testimony.
  \item[July--October 1847] Bishop de Bruillard appoints Rousselot and Orcel to investigate; they gather testimony and report.
  \item[19 September 1847] The Missionaries' chronology estimates at least 30,000 pilgrims on the first anniversary.
  \item[8 November--13 December 1847] Eight diocesan discernment conferences examine the dossier at Grenoble.
  \item[1848] Rousselot publishes a favorable report; France undergoes revolution amid continuing controversy over La Salette.
  \item[1849] The Missionaries' chronology reports 15,000 persons enrolled in the confraternity of Our Lady Reconciliatrix.
  \item[25 September 1850] Maximin meets Saint Jean-Marie Vianney at Ars; their disputed exchange produces serious doubt and later contradictory explanations.
  \item[3 July 1851] The surviving clean manuscript of Maximin's communication is dated at Grenoble; some chronologies begin the writing procedure on 2 July.
  \item[6 July 1851] Mélanie completes the extant rewritten manuscript of her communication at Grenoble.
  \item[18 July 1851] Canons Gérin and Rousselot deliver the sealed manuscripts and episcopal letter to Pius IX.
  \item[19 September 1851] Bishop Philibert de Bruillard issues the controlling positive legacy judgment and authorizes the cult of Our Lady of La Salette.
  \item[1 and 25 May 1852] De Bruillard's pastoral letter announces the shrine and Missionaries; the cornerstone is laid on 25 May.
  \item[1853--1854] Jacques-Marie-Achille Ginoulhiac succeeds de Bruillard; he declines to permit Mélanie's final profession at Corenc.
  \item[4 November 1854] Ginoulhiac issues his \emph{Instruction pastorale et Mandement}, condemning Déléon's \emph{Mémoire au Pape}, defending the established devotion and prior judgment, and separating posterior predictions from the 1846 fact.
  \item[18--19 September 1855] A near-contemporary report says Ginoulhiac described the children's mission as completed; his anniversary allocution sends pilgrims out as missionaries. The later polished shepherds-and-Church maxim summarizes this pastoral movement but is not treated as a verbatim decree.
  \item[2 February 1858] The first group of Missionaries makes religious vows. The 1879 text later claims 1858 as its authorized publication year, a statement absent from Mélanie's extant 1851 manuscript.
  \item[1865] Maximin serves six months in the Papal Zouaves; the \emph{Annales de Notre-Dame de La Salette} begins publication.
  \item[1866] Maximin publishes \emph{Ma profession de foi sur l'apparition de Notre-Dame de La Salette}.
  \item[1 March 1875] Maximin dies at Corps.
  \item[21 November 1878] Mélanie dates her long retrospective apparition narrative and expanded secret at Castellammare.
  \item[1879] The Missionaries receive Roman recognition; the basilica is consecrated and its image crowned in August. Mélanie's booklet is printed at Lecce with local permission dated 15 November.
  \item[1880] An archival Holy Office dossier, as published in later scholarship, reports measures to withdraw Mélanie's booklet and restrict similar writing or explanation; no public AAS decree or independently inspected archival original controls this entry.
  \item[14--15 December 1904] Mélanie dies at Altamura, Italy.
  \item[21 December 1915] The Holy Office issues the AAS decree against discussion of the so-called secret while expressly preserving devotion to Mary as Reconciliatrix of La Salette.
  \item[12 April and 5 June 1916] The Holy Office and Congregation of the Index proscribe \emph{La Leçon de l'Hôpital Notre-Dame d'Ypres}; the official notices print the author as ``Dr Henry Mariam'' and ``Dr Henri Mariavé.''
  \item[9--10 May 1923] The Holy Office proscribes and condemns the named 1922 Société Saint-Augustin booklet. Its AAS title prints the event date as 19 September 1845 [\emph{sic}].
  \item[8 October 1945 / 1946] Pius XII's \emph{Notre dévotion}, published in the 1946 AAS volume for the centenary, encourages the shrine, missionaries, and public call to conversion and reparation.
  \item[14 June and 15 November 1966] The CDF removes the Index's force as ecclesiastical law with attached censures while addressing conscience under natural law; it later clarifies the abrogation of old automatic book-law penalties. These acts do not by themselves decide the current force of every standalone earlier command.
  \item[1983] The current Latin Code takes effect; canon 6 abrogates Apostolic See penal laws not retained in the Code.
  \item[6 May 1996] John Paul II writes to Bishop Louis Dufaux for the 150th anniversary, interpreting La Salette as a message of hope, conversion, and reconciliation in Christ.
  \item[1999] Michel Corteville locates the sealed-era manuscripts in the historical Holy Office archive.
  \item[4 May 2000] John Paul II addresses the Missionaries' general chapter on their reconciliation charism.
  \item[2002] René Laurentin and Michel Corteville publish \emph{Découverte du secret de La Salette} with the early texts and archival material.
  \item[2016] The French Church's Nominis service reports an 18 March CDW decree for the 19 September observance; the current French national liturgical calendar directly verifies the optional celebration, while the underlying decree was not inspected here.
  \item[2021] Grenoble--Vienne and the Shrine celebrate the 175th anniversary with a diocesan Marian year and pastoral program.
  \item[19 May 2024] The DDF's new norms for alleged supernatural phenomena take effect. Their presentation explicitly cites the 1851 La Salette wording as a legacy formula; they do not reclassify the case.
  \item[4 November 2025] The DDF issues \emph{Mater Populi Fidelis}, supplying current safeguards for Marian cooperation, mediation, and intercession.
  \item[16 July 2026] DDF index, Holy See acts, Diocese, Shrine, public-message sources, and current reception rechecked. No later competent reclassification or supersession located.
\end{description}
\endgroup
